% ============================================
% BEHAVIORAL TAXONOMY & SIMULATION MODEL
% ============================================
% 👤 PERSON 3: System & Taxonomy Architect
% 
% Instructions:
% - Define the behavioral taxonomy (L1/L2 categories)
% - Explain attack stages: reconnaissance, encryption, persistence, etc.
% - Describe simulation model implementation
% - Reference Figure 3 (attack timeline/stage progression)
% - Length: 0.8-1 page
% 
% NOTE: You own Figure 3 - create attack timeline diagram
% ============================================

\section{Behavioral Taxonomy and Simulation Model}

[PERSON 3: Detail the taxonomy that drives behavioral simulation]

\subsection{Ransomware Behavioral Taxonomy}

We developed a two-level behavioral taxonomy to comprehensively model ransomware 
operations:

\textbf{Level 1 (L1): High-Level Attack Stages}
\begin{itemize}
    \item \textbf{Reconnaissance:} System profiling, file enumeration, target selection
    \item \textbf{Privilege Escalation:} Exploitation of vulnerabilities to gain elevated access
    \item \textbf{Persistence:} Installation of mechanisms to survive reboots
    \item \textbf{Encryption:} File encryption operations (the primary malicious activity)
    \item \textbf{C2 Communication:} Key exchange and ransom negotiation with attacker server
    \item \textbf{Destruction:} Shadow copy deletion, backup removal
\end{itemize}

\textbf{Level 2 (L2): Fine-Grained Behavioral Indicators}
\begin{itemize}
    \item File system operations (read, write, delete, rename patterns)
    \item Process behaviors (spawning, injection, termination)
    \item Registry modifications (persistence keys, configuration changes)
    \item Network activities (DNS queries, HTTP/HTTPS traffic, encryption protocols)
    \item Cryptographic API calls (CryptEncrypt, CryptGenKey usage patterns)
\end{itemize}

\subsection{Simulation Model}

The behavioral simulator (\texttt{simulator\_full.py}) implements a finite state 
machine progressing through L1 stages while emitting L2-level observable events. 
Each stage generates characteristic behavioral signatures that feed into the detection pipeline.

Figure~\ref{fig:timeline} illustrates typical attack stage progression and associated 
behavioral indicators.

% PERSON 3: Insert Figure 3 here
% \begin{figure}[htbp]
% \centering
% \includegraphics[width=\columnwidth]{figures/attack_timeline.pdf}
% \caption{Ransomware attack stage progression with timeline and key behavioral indicators.}
% \label{fig:timeline}
% \end{figure}

\subsection{Integration with Detection}

The \texttt{integrate\_l1\_l2.py} module bridges L1 and L2 representations, 
extracting features suitable for ML classification while preserving stage context 
for interpretability.
